\scanpage{37}
% archive sheet number: 36

\deleted{6) Soit $\Gamma=\mathcal A_5$. On a}

\[
\begin{aligned}
&\deleted{Z(\Gamma)=\{1\}}\\ 
&\deleted{\mathrm{Aut}(G,\Gamma)=N(\Gamma)=\Gamma}\\ 
&\deleted{N(\Gamma)/Z(\Gamma)=\Gamma\lhook\joinrel\longrightarrow\mathrm{Aut}(\Gamma)\simeq\unsure{\mathfrak S_5|\mathscr S_5}}
\end{aligned}
\]

Dans chacun de ces cas, on a un représentant
canonique de la catégorie fibrée envisagée, défini
sur $\mathbb Q$ : \deleted{voire $\mathbb Z$}

\uline{Cas} 1) $\mu_n$ opère sur $\mathbb G_{m\mathbb Q}$, donc sur $\widehat{\mathbb G}_{m\mathbb Q}\simeq\mathbb P^1_{\mathbb Q}$,
[l'opération s'obtient aussi en faisant opérer $\mu_{2n}$ sur
$\mathcal O_{\mathbb Q}^{2}$, \deleted{\unsure{par de|par}} la façon tautologique sur le premier
facteur $\mathcal O_{\mathbb Q}$, et \deleted{triviale} \inserted{opposée} sur le deuxième, et en
passant à $\mathbb P^1_{\mathbb Q}=\check{\mathbb P}(\mathcal O_{\mathbb Q}^2)$\unsure{$\}$|$]$}, ce qui fait passer au
quotient de $\mu_{2n}$ à $\mu_n$…]

% slanted
\margin{\uline{NB} Il n'y a de forme \deleted{non} \uline{tordue} de $\mathbb Z/n$ contenue dans un $G$ forme de $\mathrm{GP}(1)$, sur $\mathbb Q$, que si l'équation $\frac12$ cyclotomique de degré $n$ \deleted{\unsure{est de degré}} \inserted{est de degré 1}, \deleted{\unsure{sur $\mathbb Q$ ou plutôt lin.}}, ce qui signifie que $n\in\{2,3,4,6\}$. Idem pour avoir une forme non tordue de $\mathbb D_n$ dans un $G$. Pour avoir telle forme \unsure{standart|standarde|standard} $\Gamma_k\subset G$ valable pour tt $n$, on \unsure{est bien|bien} obligé de le prendre tordue !}

\uline{Cas 2, 3} $\mathbb Z/2\,\mathbin{\cdot_{1/2}}\mu_n$ opère sur $\mathbb G_{m\mathbb Q}$, donc sur
$\widehat{\mathbb G}_{m\mathbb Q}=\mathbb P^1_{\mathbb Q}$.

\margin{[\uline{ReNB} Sur un corps de base (disons) donné $k$, \uline{et $n$} donné, conditions équiv. a) $\exists\alpha\in k$ solution de l'équation $\frac12$ cyclotomique $\Psi_n(\alpha)=0$ b) $\exists$ forme non tordue de $\mathbb Z/n\subset$ forme $G$ de $\mathrm{GP}(1)$ (alors $G$ est néc. non tordue) c) Itou avec $\mathbb D_n$}

\deleted{Cas} En fait, tous ces cas sont réalisés par
les plgts naturels de $\mu_n$, $\mathbb Z/2\mathbin{\cdot_{1/2}}\mu_n$ dans
$\mathbb Z/2\mathbin{\cdot_{1/2}}\mathbb G_m$, normalisateur du tore maximal
standard de $\mathrm{GP}(1)_{\mathbb Q}\simeq\mathrm{Gl}(2)_{\mathbb Q}/\mathbb G_{m\mathbb Q}$.

\uline{Cas} 4,5 bis \inserted{\unsure{a priori|prenons}} prendre le plgmt précédent
$\mathbb D_{2\mathbb Q}=\mathbb Z/2\mathbin{\cdot_{1/2}}\mu_{2\mathbb Q}$ de $\mathrm{GP}(1)$, on prend
les normalisateurs \inserted{$\mathbb Q(i)$} de $\mathbb D_{2\mathbb Q}$ dans $\mathrm{GP}(1)$ --
celui-ci \deleted{fait} une forme \inserted{(tordue ?)} de $\mathfrak S_4$ plongée dans $\mathrm{GP}(1)$,
il faudrait regarder de plus près. [\unsure{toujours|au jugé}
\unsure{du comm. est|on comm. ait} une forme de $\mathcal A_4$,

% slanted
\margin{à choisir éventuellt un autre représentant plus commode : \deleted{\unsure{$\mathbb D$}} \unsure{le semi-dir.|le déduit} de $\mathfrak S_4$ par \unsure{$\mathcal A$|$\mathbb Z$} : mais alors $G$ n'est \unsure{pas sûr être|pas sensiblement} \unsure{bien définie|bien défini}.}
